% !TeX root = ../Book4.tex
\begin{figure}[htbp]
	\centering
	\begin{tikzpicture}[x=2.1cm,y=1.15cm,font=\small]
		\fill[GeometryBlueLight] (-.2,-1) rectangle (2.2,1);
		\draw[GeometryGray] (-.2,-1)--(2.2,-1);
		\draw[GeometryGray] (-.2,1)--(2.2,1);
		\draw[GeometryGray,dashed] (0,-1)--(0,1);
		\draw[GeometryGray,dashed] (1,-1)--(1,1);
		\draw[GeometryGray,dashed] (2,-1)--(2,1);
		\draw[GeometryBlue,thick,->] (.25,-1)--(.25,.85);
		\draw[GeometryBlue,thick] (.25,.85)--(.25,1);
		\draw[red!70!black,thick,->] (.25,-1)--(.25,-.75) .. controls (.25,-.3) and (1.25,.3) .. (1.25,.75)--(1.25,.9);
		\draw[red!70!black,thick] (1.25,.9)--(1.25,1);
		\draw[GeometryGray,->] (0,0)--(1.8,0) node[above] {\(c:\ \theta\text{ 增大}\)};
		\draw[<->] (.25,1.25)--(1.25,1.25) node[midway,above] {一个周期};
		\node[left] at (-.2,1) {\(t=1\)};
		\node[left] at (-.2,-1) {\(t=-1\)};
		\node[GeometryBlue,left] at (.25,-.55) {\(\gamma\)};
		\node[red!70!black,right] at (1.2,.55) {\(T_c\gamma\)};
		\node[below] at (0,-1) {\(0\)};
		\node[below] at (1,-1) {\(1\)};
		\node[below] at (2,-1) {\(2\)};
	\end{tikzpicture}
	\caption[Dehn 扭转的带状覆盖]{环带的带状覆盖中，蓝色横弧经正扭转变为红色弧。下端提升不变，上端提升沿核心正向增加一个周期；按 \(\theta\sim\theta+1\) 取商后，两条弧端点相同。图中定向为先向右的核心方向、后向上的横向方向。}
	\label{g4:fig:0601:dehn-strip}
\end{figure}
